% -----------------------------------------------------------------------------
% 5. Discussion and Limitations — deux implications, puis un paragraphe de
% limites dense (une proposition par limite). Budget : ~250 mots.
% -----------------------------------------------------------------------------
Two consequences follow for annotation practice. First, the reliability cost of
multi-label annotation need not be paid by abandoning multi-label annotation: it
can be recovered by consolidating the taxonomy. Second, that consolidation cannot be
delegated to the agreement signal: a purely data-driven procedure selects the
taxonomy that collapses the legal distinction the downstream
task depends on, so the criterion has to come from outside the annotation
statistics \cite{eu1993directive,micklitz2017empire}. For the evaluation of language models, the gap between the human
ceiling and the best judge (Table~\ref{tab:systems}) is large enough that judges
are not substitutable for trained annotators on this task
\cite{thalken2023legal,bavaresco2025judges}; a compact domain-specific encoder
closes much, though not all, of that gap.

\paragraph{Limitations.}
The four judges we evaluate are the four whose candidates were shown during
annotation, so the benchmark is not free of circularity \cite{choi2024llmeffect}: no
candidate was pre-filled or favoured and a third of gold labels match no proposal,
but an error shared by all four could survive, and where annotators split, the gold
coincides more often with the strongest candidate (\textsc{Fable} 0.636 against 0.607
on unanimous sentences; \textsc{Claude} 0.543 against 0.578) --- an exposure effect,
not a preference built into the protocol. Excluding \textsc{Fable} leaves
\textsc{Claude} at
$\kappa = 0.515$, which widens the gap to the human ceiling rather than narrowing
it. The layer covers one legal system and one language, and the stored judge
labels measure a historical output rather than the current state of the art. The
fine-tuned encoder is a single model with early stopping on the test fold, making
its reported agreement mildly optimistic.
Finally, the aggregated gold carries only 1.04 labels per sentence, so a
multi-label score computed on it would not measure the difficulty of the task as
annotated; the multi-label cost reported here is measured on individual votes,
where the secondary themes survive.
